\section*{Cálculo del Tiempo de Exposición en Campos de Radiación}

Para calcular el tiempo que puede permanecer una persona en un campo de radiación conocido, utilizamos la siguiente fórmula:
%
\[
t = \frac{H}{X}
\]
%
donde:
\begin{itemize}
  \item \( H \) = Dosis límite
  \item \( X \) = Radiación en el área
  \item \( t \) = Tiempo
\end{itemize}

Si queremos calcular cuánto tiempo puede permanecer un trabajador en un área con \( 30 \, \text{mr/hr} \), sin que se le permitan recibir más de \( 90 \, \text{mr} \), sustituimos los datos en la fórmula anterior y obtenemos:
%
\[
t = \frac{90 \, \text{mr}}{30 \, \text{mr/hr}} = 3 \, \text{hrs.}
\]

\section*{Cálculo de la Rapidez de Exposición a Diferentes Distancias}

Para calcular la rapidez de exposición a otra distancia diferente a la conocida utilizamos la ley del cuadrado inverso:
%
\[
X = \left(\frac{d}{D}\right)^2 \times X
\]
%
Por ejemplo, si queremos conocer la rapidez de exposición a \( 2 \, \text{m} \) de distancia de una fuente que está dando \( 100 \, \text{mr/hr} \) a \( 1 \, \text{m} \) de distancia, sustituimos los datos en la fórmula anterior y obtenemos:
%
\[
X = \left(\frac{1}{2}\right)^2 \times 100 \, \text{mr/hr} = 25 \, \text{mr/hr.}
\]

\section*{Cálculo del Tiempo Máximo de Trabajo en Distancia Determinada}

Por otra parte, si queremos calcular cuánto tiempo podemos trabajar a \( 5 \, \text{m} \) de distancia de una fuente que sabemos que está dando una rapidez de exposición de \( 2 \, \text{R/hr} \) a un metro de distancia y no queremos recibir más de \( 100 \, \text{mr} \), primero calculamos la rapidez de exposición:
%
% Calculation is pending here (to be calculated based on user's instruction using the inverse square law)
% Following portion can be completed after user has provided specific formula details or context.